% !PRINTED 319
\weberpage{319}
%<<CONT
$z = x + yi$; therein let $\varphi$ and $\psi$ be real functions of $x$ and
$y$, and now let a closed curve $f$ be drawn. We then have the theorem:

\gesperrt{XI. If the curve $f$ is divided by the curve $\varphi$ into $2m$
segments, in which $\varphi$ is alternately positive and negative, then the
number of roots of $F(z)$ lying within $f$ is at most equal to $m$;
and the same also holds when $f$ is divided by $\psi$ into $2m$ such
segments.}

The theorem is an immediate consequence of the theorem on characteristics II,
\S.\,96; for by that theorem
\[
k = \tfrac{1}{2}[E(f;\varphi,\psi) - A(f;\varphi,\psi)];
\]
and if now the curve $f$ is cut by the $\varphi$-curve in $2m$ points,
then
\[
m = \tfrac{1}{2}[E(f;\varphi,\psi) + A(f;\varphi,\psi)],
\]
hence
\[
k = m - A(f;\varphi,\psi).
\]

But here, by the theorem just mentioned, $k$ is equal to the number of roots
of $F$ lying in the interior of $f$, and $A(f;\varphi,\psi)$ is in any case a
non-negative number. In exactly the same way one can proceed with the formula
\[
k = \tfrac{1}{2}[E(f;\psi,\varphi) - A(f;\psi,\varphi)]
\]

\begin{center}
\S.\,106.

The Theorem of Rolle.
\end{center}

If the equation $f(x) = 0$ has two real roots $\alpha$ and $\beta$, of which
the first is an $a$-fold root and the second a $b$-fold root, then we can
set
\weq{1}{f(x) = (x - \alpha)^a (x - \beta)^b f_1(x)}

Now if $\alpha < \beta$, and no root of $f(x) = 0$ lies between $\alpha$ and
$\beta$, then $f(x)$ and $f_1(x)$ will have an unchanged sign throughout
this whole interval. If we form the derivative of (1), we obtain
%CONT>>
